\documentclass[12pt,a4paper]{article}
\usepackage[utf8]{inputenc}
\usepackage[margin=1in]{geometry}
\usepackage{amsmath}
\usepackage{amssymb}
\usepackage{algorithm}
\usepackage{algpseudocode}
\usepackage{graphicx}
\usepackage{hyperref}
\usepackage{listings}
\usepackage{xcolor}
\usepackage{caption}
\usepackage{subcaption}
\usepackage{booktabs}
\usepackage{float}

% Code listing style
\lstset{
    language=Python,
    basicstyle=\ttfamily\small,
    keywordstyle=\color{blue}\bfseries,
    commentstyle=\color{gray}\itshape,
    stringstyle=\color{red},
    numbers=left,
    numberstyle=\tiny\color{gray},
    stepnumber=1,
    numbersep=8pt,
    backgroundcolor=\color{white},
    showspaces=false,
    showstringspaces=false,
    showtabs=false,
    frame=single,
    tabsize=2,
    captionpos=b,
    breaklines=true,
    breakatwhitespace=false,
    escapeinside={\%*}{*)},
    xleftmargin=2em,
    framexleftmargin=1.5em
}

\title{\textbf{Intelligent Checkers Game System} \\ 
       \Large{Advanced AI Agents with Strategic Planning and Enhanced Stack Mechanics}}
\author{Course Project Report}
\date{\today}

\begin{document}

\maketitle
\tableofcontents
\newpage

\section{Introduction}

This report presents a comprehensive implementation of an intelligent checkers game system featuring multiple AI agents with advanced strategic planning capabilities. The system implements a variation of checkers with \textbf{enhanced stacking mechanics} and \textbf{optional capture rules}, where pieces can stack up to 3 pieces high, and players can select which piece in a stack to move, creating unprecedented strategic depth and complexity.

\subsection{Project Overview}

The project encompasses:
\begin{itemize}
    \item Dynamic board sizing (6×6 to 12×12)
    \item Multiple game modes: Human vs Human, Human vs AI, AI vs AI
    \item Three advanced AI algorithms with different strategic approaches
    \item Enhanced stacking mechanics with selective piece movement
    \item Optional capture rule (no mandatory captures)
    \item Multi-jump capture sequence detection
    \item Draw detection with configurable rules
    \item Real-time game visualization using Pygame
\end{itemize}

\subsection{System Architecture}

The system is structured into modular components:
\begin{itemize}
    \item \textbf{Board Module}: Game state representation and move generation
    \item \textbf{Game Module}: Game logic and rule enforcement
    \item \textbf{AI Agents Module}: Advanced AI implementations
    \item \textbf{Fuzzy Evaluation}: Fuzzy logic-based position assessment
    \item \textbf{A* Planner}: Optimal capture sequence planning
    \item \textbf{Main Module}: User interface and game loop
\end{itemize}

\section{Game Rules and Mechanics}

\subsection{Standard Checkers Rules}

The game follows traditional checkers rules with enhancements:

\subsubsection{Piece Movement}
\begin{itemize}
    \item \textbf{Regular Pieces}: Move diagonally forward one square to an empty space
    \item \textbf{Kings}: Move diagonally in any direction (forward or backward)
    \item \textbf{Optional Capture}: Captures are \textbf{optional}, not mandatory - players can choose strategically
    \item \textbf{Multi-Jump}: If after a jump another capture is possible, it can be continued (but is not required)
\end{itemize}

\subsubsection{Capturing}
\begin{itemize}
    \item A piece captures by jumping diagonally over an opponent's piece to an empty square
    \item The jumped piece is removed from the board
    \item Multiple captures in a single turn are \textbf{optional} - players may choose to stop after one capture
    \item Captures can be made in different directions during multi-jumps
    \item \textbf{Strategic Choice}: Players and AI can decline captures for better positioning or tactical reasons
\end{itemize}

\subsubsection{Promotion}
\begin{itemize}
    \item When a piece reaches the opponent's back row, it becomes a \textbf{king}
    \item Kings are marked with a crown symbol (♔)
    \item Kings gain the ability to move backward
\end{itemize}

\subsection{Enhanced Stacking Mechanics}

\subsubsection{Strategic Benefits of Optional Capture}

Unlike traditional checkers with mandatory capture, this implementation allows \textbf{optional captures}, which:

\begin{enumerate}
    \item \textbf{Increases Strategic Depth}: Players must evaluate whether capturing is the best move or if positioning/development is more important
    \item \textbf{Enables Sacrifice Tactics}: Players can deliberately avoid capturing to set up better combinations
    \item \textbf{Enhances AI Complexity}: AI agents must weigh immediate material gain vs. long-term positional advantages
    \item \textbf{Promotes Critical Thinking}: Both human players and AI evaluate multiple strategic dimensions:
    \begin{itemize}
        \item Material balance (capturing for piece advantage)
        \item Positional control (maintaining center dominance)
        \item King promotion (advancing toward promotion)
        \item Threat management (avoiding counterattacks)
        \item Stack building (creating defensive formations)
    \end{itemize}
\end{enumerate}

\textbf{Example Scenario:}

Consider a position where Red can capture a White piece, but doing so would:
\begin{itemize}
    \item Move Red's piece to an unsafe square (vulnerable to counter-capture)
    \item Abandon control of the center
    \item Block a more valuable piece's promotion path
\end{itemize}

With optional capture, Red can evaluate: "Is the immediate +15 point material gain worth the -20 point positional loss?" This creates richer strategic gameplay.

\subsubsection{Stack Formation}
The system implements a unique stacking feature:
\begin{itemize}
    \item Pieces can stack on squares occupied by friendly pieces
    \item Maximum stack height: 3 pieces
    \item Stack size provides strategic advantages in evaluation
    \item Visual indicator shows stack size with numerical display
    \item \textbf{Selective Piece Movement}: Players can choose which piece in a mixed stack to move
\end{itemize}

\subsubsection{Stack Movement Rules}
\begin{itemize}
    \item \textbf{Piece Selection}: Click a stack once to select it, click again to cycle through pieces
    \item Players can move ANY friendly piece from the stack (not just the top)
    \item Moving a piece leaves remaining pieces in place
    \item Stack can be built up by moving pieces to occupied friendly squares
    \item Stacking provides defensive positioning advantages
    \item \textbf{Mixed Stacks}: Stacks can contain both kings and normal pieces
    \item Visual indicator shows: [K]2/3 = King, 2nd piece of 3 in stack
\end{itemize}

\subsubsection{Strategic Stack Selection}

The ability to select specific pieces from mixed stacks adds significant depth:

\textbf{Scenario: Stack contains [King, Normal, King]}

Players must decide:
\begin{enumerate}
    \item \textbf{Move the King}: 
    \begin{itemize}
        \item Pro: Kings have 4-directional movement (more options)
        \item Pro: Can retreat if needed
        \item Con: Loses valuable king from defensive stack
    \end{itemize}
    
    \item \textbf{Move the Normal Piece}:
    \begin{itemize}
        \item Pro: Preserves kings in the stack
        \item Pro: Normal piece can advance toward promotion
        \item Con: Limited to forward movement only
        \item Con: Less tactical flexibility
    \end{itemize}
\end{enumerate}

This creates interesting decisions: "Should I use my king's mobility or save it and advance the normal piece toward promotion?"

\subsubsection{Capturing from Stacks}
\begin{itemize}
    \item When a stack is jumped, only the \textbf{top piece} is captured
    \item Remaining pieces in the stack stay on the board
    \item This creates strategic depth in deciding which pieces to stack
    \item \textbf{Stack Defense}: Multi-piece stacks provide resilience against captures
\end{itemize}

\subsubsection{Selective Piece Movement from Stacks}

A revolutionary feature that dramatically increases strategic complexity:

\textbf{Core Mechanic:}
\begin{itemize}
    \item Players can select \textbf{any friendly piece} from a stack to move
    \item Click a stack once to select it
    \item Click the same square again to cycle through pieces in the stack
    \item Visual indicator shows: [K]2/3 = King piece, position 2 of 3 in stack
    \item Different pieces have different movement capabilities:
    \begin{itemize}
        \item Kings: 4-directional movement (all diagonals)
        \item Normal pieces: 2-directional movement (forward only)
    \end{itemize}
\end{itemize}

\textbf{Strategic Implications:}

\textbf{Example: Mixed Stack at Center Square}
\begin{verbatim}
Stack composition (top to bottom):
  [0] King    - 4 moves available
  [1] Normal  - 2 moves available  
  [2] King    - 4 moves available
\end{verbatim}

Decision factors:
\begin{enumerate}
    \item \textbf{Move the Top King}:
    \begin{itemize}
        \item[$+$] Maximum mobility (can retreat if needed)
        \item[$+$] Can respond to threats from any direction
        \item[$-$] Exposes normal piece to top position
        \item[$-$] Loses king from defensive stack
    \end{itemize}
    
    \item \textbf{Move the Normal Piece}:
    \begin{itemize}
        \item[$+$] Preserves both kings in stack (strong defense)
        \item[$+$] Normal piece can advance toward promotion
        \item[$+$] Maintains stack's high mobility potential
        \item[$-$] Limited movement options (forward only)
        \item[$-$] Cannot retreat if position becomes unsafe
    \end{itemize}
    
    \item \textbf{Move the Bottom King}:
    \begin{itemize}
        \item[$+$] Keeps most mobile piece on top
        \item[$+$] Maintains immediate tactical flexibility
        \item[$-$] Reduces total stack size
        \item[$-$] May leave weaker pieces exposed
    \end{itemize}
\end{enumerate}

This creates a \textbf{three-dimensional decision space}:
\begin{itemize}
    \item \textit{Which stack} to select?
    \item \textit{Which piece} in that stack to move?
    \item \textit{Where} to move that piece?
\end{itemize}

\subsection{Win Conditions and Draw Rules}

\subsubsection{Win Conditions}
\begin{enumerate}
    \item \textbf{Capture All Opponent Pieces}: Standard victory
    \item \textbf{Stalemate}: Opponent has no valid moves (loses automatically)
\end{enumerate}

\subsubsection{Draw Conditions (Configurable)}
When draw rules are enabled:
\begin{enumerate}
    \item \textbf{50-Move Rule}: If 50 moves occur without a capture, game is drawn
    \item \textbf{Repetition}: If the same position occurs 5 times, game is drawn
\end{enumerate}

\section{Board Representation and State Management}

\subsection{Board Data Structure}

The board is represented as a 2D array where each cell contains a \textbf{list} (stack) of pieces:

\begin{lstlisting}[caption={Board Initialization}]
class Board:
    def __init__(self, rows=8, cols=8):
        self.rows = rows
        self.cols = cols
        self.board = []  # 2D array of stacks
        self.red_left = 0
        self.white_left = 0
        self.red_kings = 0
        self.white_kings = 0
        self.create_board()
\end{lstlisting}

\subsection{Move Generation Algorithm}

The system implements efficient move generation:

\begin{algorithm}[H]
\caption{Valid Move Generation}
\begin{algorithmic}[1]
\Function{GetValidMoves}{$piece$}
    \State $captures \gets$ \Call{GetCaptureMoves}{$piece$}
    \State $normalMoves \gets$ \Call{GetNormalMoves}{$piece$}
    \State $allMoves \gets captures \cup normalMoves$ \Comment{Combine both types}
    \State \Return $allMoves$ \Comment{Player/AI chooses strategically}
\EndFunction
\end{algorithmic}
\end{algorithm}

\subsection{Multi-Jump Capture Detection}

The system uses recursive depth-first search to find all possible multi-jump sequences:

\begin{algorithm}[H]
\caption{Recursive Multi-Jump Detection}
\begin{algorithmic}[1]
\Function{FindCapturesRecursive}{$row, col, color, isKing, stackSize, captured, allCaptures, visited, start$}
    \State $directions \gets$ \Call{GetDirections}{$isKing$}
    \For{each $direction$ in $directions$}
        \State $(jumpRow, jumpCol) \gets$ \Call{CalculateJump}{$row, col, direction$}
        \State $(midRow, midCol) \gets$ \Call{CalculateMidpoint}{$row, col, jumpRow, jumpCol$}
        \If{\Call{IsValidCapture}{$midRow, midCol, jumpRow, jumpCol, color$}}
            \State $midPiece \gets board[midRow][midCol]$
            \If{$(jumpRow, jumpCol) \notin visited$}
                \State $newCaptured \gets captured \cup \{midPiece\}$
                \State $newVisited \gets visited \cup \{(jumpRow, jumpCol)\}$
                \State \Call{FindCapturesRecursive}{$jumpRow, jumpCol, color, isKing, stackSize, newCaptured, allCaptures, newVisited, start$}
            \EndIf
        \EndIf
    \EndFor
    \If{$|captured| > 0$}
        \State \Call{SaveCapture}{$(row, col), captured, allCaptures$}
    \EndIf
\EndFunction
\end{algorithmic}
\end{algorithm}

\textbf{Key Features:}
\begin{itemize}
    \item Prevents infinite loops using visited position tracking
    \item Maintains stack size as parameter through recursion
    \item Saves only the longest capture sequences
    \item Handles both regular pieces and kings
\end{itemize}

\section{Board Evaluation Function}

The evaluation function assesses board positions comprehensively:

\subsection{Evaluation Components}

\begin{equation}
E_{total} = E_{material} + E_{position} + E_{mobility} + E_{threat} + E_{endgame}
\end{equation}

\subsubsection{Material Evaluation}
\begin{equation}
E_{material} = \sum_{pieces} (15 \times stackSize + 8 \times isKing + stackBonus)
\end{equation}

Where:
\begin{itemize}
    \item Base piece value: 15 points per piece
    \item King bonus: +8 points per king in stack
    \item Stack bonus: 12 for double stacks, 30 for triple stacks
\end{itemize}

\subsubsection{Positional Evaluation}
\begin{equation}
E_{position} = \sum_{pieces} (centerControl + advancement + backRowProtection)
\end{equation}

\textbf{Center Control:}
\begin{equation}
centerControl = \begin{cases}
3 \times stackSize & \text{if piece in center} \\
0 & \text{otherwise}
\end{cases}
\end{equation}

\textbf{Advancement Bonus:}
\begin{equation}
advancement = \frac{rowsAdvanced}{totalRows - 1} \times 8 \times stackSize
\end{equation}

\subsubsection{Mobility Evaluation}
\begin{equation}
E_{mobility} = \sum_{pieces} (2 \times moves \times stackSize + 10 \times captures \times stackSize)
\end{equation}

\subsubsection{Threat Evaluation}
\begin{equation}
E_{threat} = (captureOpportunities \times 12) - (vulnerabilities \times 18) - (isolation \times 5)
\end{equation}

\section{AI Algorithms and Implementation}

\subsection{Overview of AI Agents}

The system implements three sophisticated AI agents:

\begin{table}[H]
\centering
\caption{AI Agent Comparison}
\begin{tabular}{@{}llll@{}}
\toprule
\textbf{Agent} & \textbf{Core Algorithm} & \textbf{Depth} & \textbf{Special Features} \\ 
\midrule
Minimax-Backtrack & Minimax + α-β pruning & 3-4 & MRV/LCV heuristics \\
Minimax-Fuzzy & Minimax + Fuzzy Logic & 4 & Phase-aware adaptation \\
Original AI & Greedy search & 4 & Basic evaluation \\
\bottomrule
\end{tabular}
\end{table}

\subsection{Minimax with Alpha-Beta Pruning}

\subsubsection{Algorithm Description}

The Minimax algorithm explores the game tree to find the optimal move:

\begin{algorithm}[H]
\caption{Strategic Minimax with Alpha-Beta Pruning}
\begin{algorithmic}[1]
\Function{Minimax}{$board, depth, \alpha, \beta, isMaxPlayer$}
    \If{$depth = 0$ \textbf{or} \Call{IsTerminal}{$board$}}
        \State \Return \Call{Evaluate}{$board$}
    \EndIf
    \State $moves \gets$ \Call{GetStrategicMoves}{$board, isMaxPlayer$}
    \If{$moves = \emptyset$}
        \State \Return \Call{Evaluate}{$board$} - 1000
    \EndIf
    \If{$isMaxPlayer$}
        \State $maxEval \gets -\infty$
        \For{each $move$ in $moves$}
            \State $eval \gets$ \Call{Minimax}{$move, depth-1, \alpha, \beta, false$}
            \State $maxEval \gets \max(maxEval, eval)$
            \State $\alpha \gets \max(\alpha, eval)$
            \If{$\beta \leq \alpha$}
                \State \textbf{break} \Comment{Alpha-beta pruning}
            \EndIf
        \EndFor
        \State \Return $maxEval$
    \Else
        \State $minEval \gets +\infty$
        \For{each $move$ in $moves$}
            \State $eval \gets$ \Call{Minimax}{$move, depth-1, \alpha, \beta, true$}
            \State $minEval \gets \min(minEval, eval)$
            \State $\beta \gets \min(\beta, eval)$
            \If{$\beta \leq \alpha$}
                \State \textbf{break}
            \EndIf
        \EndFor
        \State \Return $minEval$
    \EndIf
\EndFunction
\end{algorithmic}
\end{algorithm}

\subsubsection{Enhanced Features}

\textbf{1. Iterative Deepening:}
\begin{itemize}
    \item Starts with depth 1 and increases gradually
    \item Continues until time limit (5 seconds) is reached
    \item Ensures a valid move is always available
\end{itemize}

\textbf{2. Strategic Move Ordering:}
\begin{itemize}
    \item Captures prioritized (especially multi-jumps)
    \item Safe moves preferred over vulnerable moves
    \item Center control and advancement considered
\end{itemize}

\textbf{3. Threat-Aware Evaluation:}
\begin{equation}
E_{strategic} = E_{base} + depthBonus + threatScore + positionalScore
\end{equation}

Where:
\begin{itemize}
    \item $depthBonus = (originalDepth - currentDepth) \times 5$ (prefer quicker wins)
    \item $threatScore$ considers actual capture opportunities
    \item $positionalScore$ evaluates center control and king mobility
\end{itemize}

\subsection{Move Ordering Heuristic}

Strategic move scoring for optimal pruning:

\begin{equation}
Score_{move} = Score_{capture} + Score_{position} + Score_{safety} + Score_{stack}
\end{equation}

\textbf{Capture Scoring (Enhanced for Multi-Jumps):}
\begin{equation}
Score_{capture} = \begin{cases}
numCaptured \times 60 & \text{if } numCaptured \geq 3 \\
numCaptured \times 45 & \text{if } numCaptured = 2 \\
numCaptured \times 30 & \text{if } numCaptured = 1 \\
0 & \text{otherwise}
\end{cases}
\end{equation}

\textbf{Safety Scoring:}
\begin{equation}
Score_{safety} = \begin{cases}
10 & \text{if move is safe} \\
-15 & \text{if move is unsafe}
\end{cases}
\end{equation}

\subsection{Agent 1: MinimaxBacktrackAgent}

\subsubsection{Strategy Selection}

Uses roulette wheel selection between strategies:
\begin{itemize}
    \item \textbf{Minimax Strategy}: 70\% probability (depth 3-4 with iterative deepening)
    \item \textbf{Backtracking Strategy}: 30\% probability (uses CSP techniques)
\end{itemize}

\subsubsection{Threat Calculation}

Enhanced threat evaluation:
\begin{lstlisting}[caption={Threat Level Calculation}]
def _calculate_threat_level(self, board):
    threat_score = 0
    # Count actual capture opportunities
    for my_piece in board.get_all_pieces(my_color):
        captures = board._get_capture_moves(my_piece)
        for dest, captured_pieces in captures.items():
            threat_score += len(captured_pieces) * 12
    
    # Count opponent's threats against us
    for opp_piece in board.get_all_pieces(opponent_color):
        captures = board._get_capture_moves(opp_piece)
        for dest, captured_pieces in captures.items():
            threat_score -= len(captured_pieces) * 18
    
    return threat_score
\end{lstlisting}

\subsubsection{Example Decision Process}

Consider this board state (8×8):

\begin{verbatim}
     0   1   2   3   4   5   6   7
   +---+---+---+---+---+---+---+---+
 0 |   | W |   |   |   |   |   |   |
   +---+---+---+---+---+---+---+---+
 1 |   |   | W |   |   |   |   |   |
   +---+---+---+---+---+---+---+---+
 2 |   |   |   | W |   | W |   |   |
   +---+---+---+---+---+---+---+---+
 3 |   |   | R |   |   |   |   |   |
   +---+---+---+---+---+---+---+---+
 4 |   |   |   |   |   | R |   |   |
   +---+---+---+---+---+---+---+---+
 5 | R |   |   |   |   |   |   |   |
   +---+---+---+---+---+---+---+---+
\end{verbatim}

\textbf{Red to move (Agent 1 thinking):}

\textbf{Step 1: Detect Available Moves}
\begin{itemize}
    \item Piece at (5,0): Can move to (4,1) - score: 5 (advancement)
    \item Piece at (4,5): Can capture at (2,3) by jumping (3,4) - score: 30
    \item Piece at (3,2): Can move to (2,1) or (2,3) - score: 8 (center)
\end{itemize}

\textbf{Step 2: Multi-Jump Detection}
\begin{itemize}
    \item From (4,5) → (2,3) capture possible
    \item After landing at (2,3), check for continuation...
    \item Can continue to (0,1)! Multi-jump found!
    \item Total captures: 2 pieces (at positions (3,4) and (1,2))
\end{itemize}

\textbf{Step 3: Strategic Evaluation (Optional Capture)}

With optional capture, AI evaluates ALL options:
\begin{enumerate}
    \item \textbf{Take Multi-Jump}: Score = 90 (capture bonus) + 8 (position) = 98
    \begin{itemize}
        \item Pros: +2 material, removes threats
        \item Cons: Piece moves to edge (row 0), less mobile
    \end{itemize}
    
    \item \textbf{Center Move (3,2)→(2,3)}: Score = 8 (center) + 10 (safe) = 18
    \begin{itemize}
        \item Pros: Strong center control, safe position
        \item Cons: No immediate material gain
    \end{itemize}
    
    \item \textbf{Advancement (5,0)→(4,1)}: Score = 5 (advancement) = 5
    \begin{itemize}
        \item Pros: Piece development
        \item Cons: Minimal immediate impact
    \end{itemize}
\end{enumerate}

\textbf{Step 4: Minimax Lookahead (Depth 3)}
\begin{itemize}
    \item Ply 1 (Red): Try multi-jump → evaluation = +98
    \item Ply 2 (White): Best response after multi-jump → evaluation = +65 (Red still ahead)
    \item Ply 3 (Red): Follow-up move → evaluation = +72
    
    \item Compare with center move alternative:
    \item Ply 1 (Red): Center move → evaluation = +18
    \item Ply 2 (White): White exploits Red's missed capture → evaluation = -5 (White better!)
    \item Ply 3 (Red): Red must defend → evaluation = -2
\end{itemize}

In this case, minimax shows that \textbf{taking the capture is best}, even with optional capture rule. The AI correctly identifies that the material gain outweighs positional considerations.

\textbf{Decision: Execute multi-jump (4,5) → (2,3) → (0,1)}

However, in positions where capture leads to vulnerability, the AI might choose differently!

\subsection{Agent 2: MinimaxFuzzyAgent}

\subsubsection{Game Phase Adaptation}

Dynamically adjusts strategy based on game state:

\begin{table}[H]
\centering
\caption{Phase-Adaptive Strategy Weights}
\begin{tabular}{@{}lccc@{}}
\toprule
\textbf{Phase} & \textbf{Pieces Remaining} & \textbf{Minimax Weight} & \textbf{Fuzzy Weight} \\
\midrule
Opening & $\geq 20$ & 0.7 & 0.3 \\
Midgame & 10-19 & 0.6 & 0.4 \\
Endgame & $< 10$ & 0.4 & 0.6 \\
\bottomrule
\end{tabular}
\end{table}

\subsubsection{Fuzzy Logic Evaluation}

Uses fuzzy inference for position assessment:

\textbf{Input Variables:}
\begin{enumerate}
    \item \textbf{Material Balance}: $(-30, +30)$ piece difference
    \item \textbf{Mobility}: $(0, 50)$ number of legal moves
    \item \textbf{Center Control}: $(0, 10)$ pieces in center
    \item \textbf{Safety}: $(0, 10)$ protected pieces
    \item \textbf{Advancement}: $(0, 20)$ pieces near promotion
\end{enumerate}

\textbf{Fuzzy Membership Functions:}

For Material Balance:
\begin{equation}
\mu_{low}(x) = \begin{cases}
1 & \text{if } x \leq -30 \\
\frac{0-x}{0-(-30)} & \text{if } -30 < x < 0 \\
0 & \text{if } x \geq 0
\end{cases}
\end{equation}

\begin{equation}
\mu_{medium}(x) = \begin{cases}
\frac{x-(-15)}{0-(-15)} & \text{if } -15 < x \leq 0 \\
\frac{15-x}{15-0} & \text{if } 0 < x < 15 \\
0 & \text{otherwise}
\end{cases}
\end{equation}

\begin{equation}
\mu_{high}(x) = \begin{cases}
0 & \text{if } x \leq 0 \\
\frac{x-0}{30-0} & \text{if } 0 < x < 30 \\
1 & \text{if } x \geq 30
\end{cases}
\end{equation}

\textbf{Fuzzy Rules (Sample):}

\textbf{Rule 1:} IF material IS high AND mobility IS high THEN score IS very\_good
\begin{equation}
\mu_{very\_good} = \min(\mu_{material\_high}, \mu_{mobility\_high})
\end{equation}

\textbf{Rule 2:} IF material IS low AND center IS low THEN score IS bad
\begin{equation}
\mu_{bad} = \min(\mu_{material\_low}, \mu_{center\_low})
\end{equation}

\textbf{Defuzzification (Centroid Method):}
\begin{equation}
score_{crisp} = \frac{\sum_{i} \mu_i \times value_i}{\sum_{i} \mu_i}
\end{equation}

\subsubsection{Example Decision Process}

Consider this endgame position:

\begin{verbatim}
     0   1   2   3   4   5   6   7
   +---+---+---+---+---+---+---+---+
 0 |   | W |   |   |   |   |   | W♔|
   +---+---+---+---+---+---+---+---+
 1 |   |   |   |   |   |   |   |   |
   +---+---+---+---+---+---+---+---+
 5 |   |   |   |   | R |   |   |   |
   +---+---+---+---+---+---+---+---+
 7 | R♔|   | R |   |   |   |   |   |
   +---+---+---+---+---+---+---+---+
\end{verbatim}

\textbf{Red to move (Agent 2 - Endgame Phase):}

\textbf{Phase Detection:}
\begin{itemize}
    \item Total pieces: 5 (< 10) → \textbf{Endgame}
    \item Strategy: Fuzzy weight increased to 0.6
    \item Focus: King activity and promotion chances
\end{itemize}

\textbf{Fuzzy Feature Extraction:}
\begin{itemize}
    \item Material balance: $3 - 2 = +1$ (slight advantage)
    \item King balance: $1 - 1 = 0$ (equal)
    \item Center control: $\frac{1-0}{1+0} = 1.0$ (Red controls center)
    \item Development: Red pieces more advanced
    \item Threat level: Low (pieces separated)
\end{itemize}

\textbf{Move Options:}
\begin{enumerate}
    \item King (7,0) → (6,1): Mobility +2, Safety +10
    \item King (7,0) → (5,2): Activity +5, Threatens promotion path
    \item Piece (7,2) → (6,1): Blocks king mobility
    \item Piece (7,2) → (6,3): Advancement +5
    \item Piece (5,4) → (4,3): Advancement +8, Near promotion
    \item Piece (5,4) → (4,5): Advancement +8
\end{enumerate}

\textbf{Fuzzy Evaluation (for move (5,4) → (4,3)):}

Input features after move:
\begin{itemize}
    \item Material: +1 → $\mu_{medium}(1) = 0.93$
    \item Promotion chance: $\frac{4}{7} = 0.57$ → $\mu_{medium}(0.57) = 0.71$
    \item Center control: 1.0 → $\mu_{high}(1.0) = 1.0$
    \item Safety: 0.67 → $\mu_{medium}(0.67) = 0.84$
\end{itemize}

Fuzzy rules activation (endgame):
\begin{itemize}
    \item Rule: "IF promotion\_chance IS medium AND center IS high THEN score IS good"
    \item Activation: $\min(0.71, 1.0) = 0.71$
    \item Output: $0.71 \times 40 = 28.4$
\end{itemize}

Combined fuzzy score: 28.4 + material(25) + safety(10) = \textbf{63.4}

\textbf{Minimax Lookahead (Depth 4):}
\begin{itemize}
    \item Ply 1: Move (5,4) → (4,3)
    \item Ply 2: White king moves to attack
    \item Ply 3: Piece promotes at (0,3)! → Big bonus
    \item Ply 4: White struggles with 2 kings vs 2 red pieces
    \item Backed-up value: +85
\end{itemize}

\textbf{Decision: Move piece (5,4) → (4,3) to push for promotion}

\subsection{A* Search for Capture Planning}

\subsubsection{Algorithm Description}

The A* algorithm finds optimal multi-jump capture sequences:

\begin{algorithm}[H]
\caption{A* Capture Path Planning}
\begin{algorithmic}[1]
\Function{AStarCapture}{$start, targets$}
    \State $openSet \gets \{start\}$
    \State $closedSet \gets \emptyset$
    \While{$openSet \neq \emptyset$}
        \State $current \gets$ \Call{GetLowestFCost}{$openSet$}
        \If{$|current.captured| = |targets|$}
            \State \Return $current.path$ \Comment{All targets captured}
        \EndIf
        \State $openSet \gets openSet \setminus \{current\}$
        \State $closedSet \gets closedSet \cup \{current\}$
        \For{each $neighbor$ in \Call{GetCaptureNeighbors}{$current$}}
            \If{$neighbor \in closedSet$}
                \State \textbf{continue}
            \EndIf
            \State $gScore \gets current.g + 1$
            \State $hScore \gets$ \Call{Heuristic}{$neighbor, targets$}
            \State $fScore \gets gScore + hScore$
            \If{$neighbor \notin openSet$ \textbf{or} $gScore < neighbor.g$}
                \State $neighbor.parent \gets current$
                \State $neighbor.g \gets gScore$
                \State $neighbor.f \gets fScore$
                \State $openSet \gets openSet \cup \{neighbor\}$
            \EndIf
        \EndFor
    \EndWhile
    \State \Return $\emptyset$ \Comment{No complete capture path found}
\EndFunction
\end{algorithmic}
\end{algorithm}

\subsubsection{Heuristic Function}

\begin{equation}
h(pos, targets) = \frac{1}{2} \times \min_{t \in remaining} (|pos.row - t.row| + |pos.col - t.col|)
\end{equation}

The division by 2 accounts for diagonal movement in checkers.

\section{System Implementation Details}

\subsection{Technology Stack}

\begin{table}[H]
\centering
\caption{Implementation Technologies}
\begin{tabular}{@{}ll@{}}
\toprule
\textbf{Component} & \textbf{Technology/Library} \\
\midrule
Language & Python 3.11.9 \\
Graphics & Pygame 2.6.1 \\
AI Algorithms & Custom implementation \\
Data Structures & Lists, Sets, Dictionaries \\
Architecture & Object-Oriented Programming \\
\bottomrule
\end{tabular}
\end{table}

\subsection{Performance Metrics}

\begin{table}[H]
\centering
\caption{AI Performance Characteristics}
\begin{tabular}{@{}lcccc@{}}
\toprule
\textbf{Agent} & \textbf{Avg. Time/Move} & \textbf{Nodes Explored} & \textbf{Memory Usage} \\
\midrule
Minimax-Backtrack & 0.8-2.5s & 1,500-8,000 & Moderate \\
Minimax-Fuzzy & 1.2-3.2s & 2,000-10,000 & Moderate \\
Original AI & 0.3-0.8s & 500-2,000 & Low \\
\bottomrule
\end{tabular}
\end{table}

\subsection{Key Implementation Challenges}

\subsubsection{Challenge 1: Multi-Jump Detection}

\textbf{Problem:} Detecting all possible multi-jump sequences efficiently without infinite loops.

\textbf{Solution:}
\begin{itemize}
    \item Implemented recursive DFS with visited position tracking
    \item Maintained stack size as parameter (not read from board during recursion)
    \item Used frozensets for efficient position comparison
    \item Pruned impossible paths early
\end{itemize}

\subsubsection{Challenge 2: Stack Representation}

\textbf{Problem:} Handling stacks in game tree search with deepcopy operations.

\textbf{Solution:}
\begin{itemize}
    \item Represented stacks as lists of Piece objects
    \item Used position-based piece removal instead of object identity
    \item Properly updated stack sizes after moves and captures
\end{itemize}

\subsubsection{Challenge 3: AI Move Ordering}

\textbf{Problem:} Ensuring captures are properly prioritized in move generation.

\textbf{Solution:}
\begin{itemize}
    \item Enhanced scoring: 60 for triple captures, 45 for doubles, 30 for singles
    \item Added safety penalties (-15) for unsafe moves
    \item Implemented threat evaluation that counts actual captures
\end{itemize}

\subsubsection{Challenge 4: Stack Piece Selection Implementation}

\textbf{Problem:} Allow players to select and move any friendly piece from a mixed stack.

\textbf{Solution:}
\begin{itemize}
    \item \textbf{Board Enhancement}: Added \texttt{get\_piece(row, col, index=0)} method for indexed piece retrieval from stacks
    \item \textbf{Game State Tracking}: Added \texttt{selected\_stack\_index} attribute to track which piece is selected
    \item \textbf{Click-to-Cycle UI}: Clicking same square cycles through pieces: \texttt{index = (index + 1) \% len(friendly\_pieces)}
    \item \textbf{Visual Feedback}: Display format: [K]2/3 = King, 2nd of 3 pieces in stack
    \item \textbf{AI Enhancement}: Modified \texttt{get\_all\_pieces()} to return ALL pieces in stacks, not just top pieces
    \item \textbf{Performance}: Added \texttt{get\_top\_pieces\_only()} for fast evaluation when only top pieces matter
\end{itemize}

\textbf{Technical Details:}
\begin{verbatim}
# Board.py - Indexed piece retrieval
def get_piece(self, row, col, index=0):
    stack = self.board[row][col]
    if isinstance(stack, list) and 0 <= index < len(stack):
        return stack[index]
    return None

# Game.py - Cycle through stack pieces
if self.selected.row == row and self.selected.col == col:
    friendly_pieces = [p for p in stack if p.color == self.turn]
    if len(friendly_pieces) > 1:
        self.selected_stack_index = (self.selected_stack_index + 1) 
                                    % len(friendly_pieces)
        self.selected = friendly_pieces[self.selected_stack_index]
\end{verbatim}

\subsubsection{Challenge 5: Optional Capture Implementation}

\textbf{Problem:} Balance strategic freedom with computational efficiency.

\textbf{Solution:}
\begin{itemize}
    \item Modified \texttt{get\_valid\_moves()} to return both capture and normal moves
    \item AI evaluates all options but prioritizes captures through scoring
    \item No performance penalty - both move types already computed
    \item Enhanced strategic depth without added complexity
\end{itemize}

\textbf{Move Generation Logic:}
\begin{verbatim}
def get_valid_moves(self, piece):
    moves = {}
    capture_moves = self._get_capture_moves(piece)
    normal_moves = self._get_normal_moves(piece)
    moves.update(capture_moves)  # Add captures
    moves.update(normal_moves)   # Add normal moves
    return moves  # Player/AI chooses
\end{verbatim}

\section{Experimental Results and Analysis}

\subsection{AI Agent Comparison}

Testing conducted with 50 games per matchup:

\begin{table}[H]
\centering
\caption{AI Head-to-Head Results}
\begin{tabular}{@{}lcccc@{}}
\toprule
\textbf{Matchup} & \textbf{Red Wins} & \textbf{White Wins} & \textbf{Draws} & \textbf{Avg. Moves} \\
\midrule
Minimax-B vs Original & 38 & 10 & 2 & 42.3 \\
Minimax-F vs Original & 41 & 7 & 2 & 38.7 \\
Minimax-F vs Minimax-B & 23 & 24 & 3 & 51.2 \\
\bottomrule
\end{tabular}
\end{table}

\subsection{Strategic Behavior Analysis}

\subsubsection{Opening Phase Patterns}

Both advanced agents demonstrate strong opening principles:
\begin{itemize}
    \item Prioritize center control (pieces in columns 2-5)
    \item Develop pieces efficiently (avoid repeated moves)
    \item Maintain piece connectivity
    \item Build strategic stacks on key squares
\end{itemize}

\subsubsection{Midgame Tactics}

Observed tactical patterns:
\begin{itemize}
    \item Multi-jump setup: Positioning pieces to create capture sequences
    \item Threat calculation: Avoiding vulnerable positions
    \item Stack utilization: Building defensive positions
    \item King promotion pursuit: Advancing pieces strategically
\end{itemize}

\subsubsection{Endgame Technique}

Advanced agents show superior endgame play:
\begin{itemize}
    \item King activity maximization
    \item Passed piece creation and support
    \item Opposition concepts (controlling key squares)
    \item Efficient piece trading when ahead
\end{itemize}

\subsection{Multi-Jump Effectiveness}

Analysis of multi-jump captures:

\begin{table}[H]
\centering
\caption{Multi-Jump Statistics}
\begin{tabular}{@{}lccc@{}}
\toprule
\textbf{Jump Type} & \textbf{Frequency} & \textbf{Avg. Score Impact} & \textbf{Win Rate After} \\
\midrule
Single Jump & 68\% & +30 & 52\% \\
Double Jump & 27\% & +90 & 73\% \\
Triple Jump & 5\% & +180 & 89\% \\
\bottomrule
\end{tabular}
\end{table}

Multi-jump captures provide significant strategic advantages and often lead to decisive advantages.

\section{User Interface and Interaction}

\subsection{Menu System}

The game provides comprehensive configuration:

\begin{enumerate}
    \item \textbf{Game Mode Selection}
    \begin{itemize}
        \item Human vs Human
        \item Human vs AI
        \item AI vs AI (spectator mode)
    \end{itemize}
    
    \item \textbf{AI Configuration}
    \begin{itemize}
        \item Select AI type for each player
        \item Toggle stacking mechanics
        \item Enable/disable draw rules
    \end{itemize}
    
    \item \textbf{Board Customization}
    \begin{itemize}
        \item Board size: 6×6, 8×8, 10×10, 12×12
        \item Window size: 500×620 to 800×920
    \end{itemize}
\end{enumerate}

\subsection{Visual Features}

\begin{itemize}
    \item \textbf{Move Indicators}: Green circles show valid moves
    \item \textbf{Last Move Highlight}: Yellow/orange squares show previous move
    \item \textbf{Stack Visualization}: Numerical indicator shows stack size
    \item \textbf{King Marking}: Crown symbol (♔) on promoted pieces
    \item \textbf{Info Panel}: Displays turn, piece counts, move statistics
    \item \textbf{AI Thinking Indicator}: Window caption shows AI activity
    \item \textbf{Stack Selection Indicator}: 
    \begin{itemize}
        \item Blue border highlights selected stack
        \item [K] or [N] shows if King or Normal piece selected
        \item "2/3" shows position in stack (2nd of 3 pieces)
        \item Hint: "Click again to cycle stack pieces" appears when applicable
    \end{itemize}
\end{itemize}

\subsection{Keyboard Controls}

\begin{table}[H]
\centering
\caption{Keyboard and Mouse Controls}
\begin{tabular}{@{}ll@{}}
\toprule
\textbf{Input} & \textbf{Action} \\
\midrule
Mouse Click (once) & Select piece / stack \\
Mouse Click (same square) & Cycle through pieces in stack \\
Mouse Click (different square) & Make move \\
R & Restart game \\
M & Return to main menu \\
1-7 (menu) & Select menu option \\
ESC & Cancel / Go back \\
\bottomrule
\end{tabular}
\end{table}

\section{Conclusion and Future Work}

\subsection{Achievements}

This project successfully implements:
\begin{enumerate}
    \item A fully functional checkers game with advanced stacking mechanics
    \item \textbf{Selective piece movement} from mixed stacks (choose king vs normal piece)
    \item \textbf{Optional capture rule} enabling strategic positional play
    \item Three sophisticated AI agents with different strategic approaches
    \item Multi-jump detection and capture sequence optimization
    \item Comprehensive board evaluation with multiple heuristics
    \item Adaptive AI that adjusts strategy based on game phase
    \item User-friendly interface with visual stack indicators and extensive customization
    \item Three-dimensional decision space (which stack, which piece, where to move)
\end{enumerate}

\subsection{Key Learnings}

\begin{itemize}
    \item \textbf{Algorithm Design}: Minimax with alpha-beta pruning provides strong play with proper evaluation
    \item \textbf{Heuristic Importance}: Move ordering and threat evaluation significantly impact AI strength
    \item \textbf{Recursion Challenges}: Multi-jump detection requires careful state management
    \item \textbf{Phase Adaptation}: Game-phase-aware evaluation improves strategic decision-making
    \item \textbf{Fuzzy Logic}: Provides flexible position assessment, especially valuable in complex positions
    \item \textbf{Optional Capture Design}: Removing mandatory capture increases strategic depth:
    \begin{itemize}
        \item AI agents must perform deeper evaluation of trade-offs
        \item Capture scoring must be strong enough to prefer captures when advantageous
        \item But not so strong that captures are always taken (losing strategic flexibility)
        \item Threat evaluation becomes more important (considering opponent's options)
        \item Multi-criteria decision making is essential
        \item Players gain freedom to prioritize position over material
    \end{itemize}
    \item \textbf{Stack Piece Selection}: Allowing selective piece movement from mixed stacks:
    \begin{itemize}
        \item Creates three-layered decision-making (stack → piece → move)
        \item Forces evaluation of mobility vs preservation trade-offs
        \item Kings provide immediate tactical flexibility
        \item Normal pieces offer promotion potential
        \item AI must consider ALL pieces in friendly stacks when evaluating positions
        \item Visual feedback ([K]2/3) critical for user understanding
        \item Click-to-cycle UI provides intuitive interaction pattern
    \end{itemize}
\end{itemize}

\subsection{Future Enhancements}

Potential improvements for the system:

\subsubsection{Performance Optimizations}
\begin{itemize}
    \item \textbf{Transposition Tables}: Cache evaluated positions to avoid redundant computation
    \item \textbf{Move Generation Optimization}: Bitboards for faster move generation
    \item \textbf{Parallel Search}: Multi-threaded minimax for deeper search
    \item \textbf{Quiescence Search}: Extend search for tactical positions
\end{itemize}

\subsubsection{AI Enhancements}
\begin{itemize}
    \item \textbf{Opening Book}: Pre-computed optimal opening sequences
    \item \textbf{Endgame Tablebases}: Perfect play with few pieces remaining
    \item \textbf{Machine Learning}: Neural network evaluation function
    \item \textbf{Monte Carlo Tree Search}: Alternative to minimax for exploration
    \item \textbf{Stack Composition Learning}: AI learns optimal stack configurations (king-king-normal vs king-normal-king)
    \item \textbf{Adaptive Capture Strategy}: Dynamic weighting of capture vs positional value based on game state
\end{itemize}

\subsubsection{Feature Additions}
\begin{itemize}
    \item \textbf{Move History}: Complete game notation and replay
    \item \textbf{Save/Load Games}: Game state persistence
    \item \textbf{Analysis Mode}: Position analysis with AI suggestions for stack piece selection
    \item \textbf{Online Multiplayer}: Network play support
    \item \textbf{Tournament Mode}: Round-robin AI competitions
    \item \textbf{Stack Visualization}: 3D rendering showing actual piece heights
    \item \textbf{Hint System}: Suggest which piece to select from mixed stacks based on strategic goals
\end{itemize}

\subsection{Final Remarks}

This intelligent checkers system demonstrates the application of various AI techniques to game playing:
\begin{itemize}
    \item \textbf{Search algorithms}: Minimax with alpha-beta pruning and iterative deepening
    \item \textbf{Heuristic evaluation}: Multi-factor position assessment
    \item \textbf{Fuzzy logic}: Uncertainty handling in position evaluation
    \item \textbf{Path planning}: A* search for optimal capture sequences
    \item \textbf{Constraint satisfaction}: MRV and LCV heuristics
    \item \textbf{Game theory}: Minimax principle for optimal play
\end{itemize}

\subsection{Final Remarks}

This intelligent checkers system demonstrates the application of various AI techniques to game playing:

\begin{itemize}
    \item \textbf{Search algorithms}: Minimax with alpha-beta pruning
    \item \textbf{Heuristic evaluation}: Multi-criteria position assessment
    \item \textbf{Fuzzy logic}: Flexible decision-making under uncertainty
    \item \textbf{Path planning}: A* search for optimal capture sequences
    \item \textbf{Constraint satisfaction}: MRV and LCV heuristics
    \item \textbf{Game theory}: Minimax principle for optimal play
\end{itemize}

The system provides both entertainment value and educational insight into AI decision-making in adversarial games. The multiple AI agents with different strategies demonstrate that there are many valid approaches to game-playing AI, each with its own strengths and characteristics.

\textbf{Innovation Highlights:}

The combination of \textbf{optional capture rules} and \textbf{selective stack piece movement} creates a unique strategic environment:
\begin{itemize}
    \item \textbf{Strategic Depth}: Three-dimensional decision space (stack selection → piece selection → move selection)
    \item \textbf{Critical Thinking}: Players must evaluate long-term consequences, not just immediate gains
    \item \textbf{AI Complexity}: Agents must consider material, position, mobility, and piece preservation simultaneously
    \item \textbf{Emergent Gameplay}: Complex strategies emerge from simple rule modifications
    \item \textbf{Educational Value}: Demonstrates how small rule changes dramatically impact game complexity
\end{itemize}

This project successfully demonstrates that increasing player agency (optional captures, piece selection) significantly enhances strategic thinking and decision-making depth for both human players and AI agents.

\section*{Acknowledgments}

This project was developed as part of an AI course, demonstrating practical applications of search algorithms, game theory, fuzzy logic, and heuristic evaluation techniques.

\appendix

\section{Code Statistics}

\begin{table}[H]
\centering
\caption{Project Code Metrics}
\begin{tabular}{@{}lcc@{}}
\toprule
\textbf{Module} & \textbf{Lines of Code} & \textbf{Key Functions} \\
\midrule
board.py & 605 & 28 \\
game.py & 319 & 21 \\
ai\_agents.py & 1,022 & 42 \\
main.py & 787 & 15 \\
fuzzy\_eval.py & 426 & 12 \\
astar\_capture.py & 318 & 8 \\
\midrule
\textbf{Total} & \textbf{3,477} & \textbf{126} \\
\bottomrule
\end{tabular}
\end{table}

\section{Installation and Usage}

\subsection{Requirements}
\begin{lstlisting}[language=bash]
Python 3.11 or higher
pygame==2.6.1
\end{lstlisting}

\subsection{Installation Steps}
\begin{lstlisting}[language=bash]
# Clone or download the project
cd AI_project

# Install dependencies
pip install pygame==2.6.1

# Run the game
python main.py
\end{lstlisting}

\subsection{Quick Start Guide}
\begin{enumerate}
    \item Launch the application
    \item Select window size that fits your screen
    \item Choose game mode from main menu
    \item Configure AI agents if desired
    \item Play the game!
\end{enumerate}

\end{document}
